% =========================================================
% Proof Stub: Absolute Convergence is Rearrangement Invariant
% Source: volume-iii/analysis/sequences/notes/series/notes-absolute-convergence.tex
% =========================================================

% *** HARDER PROOF — attempt after foundational results settled ***

\clearpage
\phantomsection
\hypertarget{proof-RA-ABB-T-abs-conv-rearrangement-inv}{}

\subsubsection[Absolute Convergence is Rearrangement Invariant]{Proof --- Absolute Convergence is Rearrangement Invariant}

\bigskip

\noindent
\textbf{Source.} See \texttt{notes-absolute-convergence.tex}.

\vspace{0.75em}

\noindent
\textbf{Goal.} If $\sum a_n$ converges absolutely and $\sigma$ is any bijection $\mathbb{N} \to \mathbb{N}$, then $\sum a_{\sigma(n)}$ converges to the same sum.

\vspace{0.75em}

\noindent
\textbf{Logical form.}
\[
  \sum |a_n| < \infty \Rightarrow \sum a_{\sigma(n)} = \sum a_n \text{ for every bijection } \sigma.
\]

\vspace{0.75em}

\noindent
\textbf{Proof strategy.} Given $\varepsilon > 0$, find $N$ past which the tail of $\sum |a_n|$ is $< \varepsilon$. Show partial sums of the rearrangement, once past a large enough index, include all of $a_1, \ldots, a_N$ plus a bounded tail.

\vspace{0.75em}

\noindent
\textbf{Proof.}
\begin{proof}
\Claim If $\sum a_n$ converges absolutely and $\sigma$ is any bijection $\mathbb{N} \to \mathbb{N}$, then $\sum a_{\sigma(n)}$ converges to the same sum.

\Given 

\Goal 

\Strategy Given $\varepsilon > 0$, find $N$ past which the tail of $\sum |a_n|$ is $< \varepsilon$. Show partial sums of the rearrangement, once past a large enough index, include all of $a_1, \ldots, a_N$ plus a bounded tail.

\medskip

% --- get N: ∑_{n>N} |a_n| < ε ---
% --- let M be such that {1,...,N} ⊆ σ({1,...,M}) ---
% --- for m ≥ M: ∑_{n=1}^m a_{σ(n)} includes all of a_1,...,a_N ---
% --- difference from ∑_{n=1}^N a_n is bounded by ∑_{n>N} |a_n| < ε ---
% --- both partial sum sequences are Cauchy and share accumulation points; argue they share the limit ---

\end{proof}

\begin{remark*}[Proof shape]
The absolute convergence of the tail ensures rearrangements can only pick up an $\varepsilon$-small error from the tail.
\end{remark*}

\begin{remark*}[Dependencies]
\begin{itemize}
  \item Absolute convergence: $\sum |a_n| < \infty$.
  \item Cauchy Criterion for series.
  \item $\{1,\ldots,N\} \subseteq \sigma(\{1,\ldots,M\})$ for some finite $M$ (since $\sigma$ is a bijection).
\end{itemize}
\end{remark*}
